\documentclass[11pt]{article}
\usepackage{fontspec,amsmath,amssymb,amsthm,hyperref}
\usepackage[margin=1in]{geometry}
\setlength{\emergencystretch}{3em}
\title{The actual L\ensuremath{{}^{2}} Hamiltonian comparator has the exact Gibbs law}
\author{LangevinLean mathematical exposition}
\date{}
\begin{document}\maketitle
\noindent Source module: \texttt{KineticLangevin/LpHamiltonianLaw.lean}. Compile this proof document with XeLaTeX or LuaLaTeX. The accompanying HTML pairs the mathematical argument with the exact Lean statements and records the scope of the formalization.


This module closes the probabilistic reference obligations for the Hamiltonian part of the OBABO bias argument. It constructs the actual L\ensuremath{{}^{2}} Picard flow, identifies its finite-dimensional output law, derives exact Gibbs stationarity and velocity/score moments, and applies the proved local O–Verlet–O comparison. No stationary reference-law or local-error bound is postulated.


\section{Identify an L\ensuremath{{}^{2}} limit with a pointwise numerical limit}


Suppose actual L\ensuremath{{}^{2}} phase variables \(S_n\) converge to S, and their representatives agree almost everywhere with measurable vector approximants \(f_n\) converging almost everywhere to Y. For any continuous linear functional \(\ell\), the characteristic test \(e^{i\ell(z)}\) has norm one. Bounded dominated convergence identifies the characteristic-function limit of the pointwise approximants, while the proved L\ensuremath{{}^{2}} characteristic continuity identifies the limit of the L\ensuremath{{}^{2}} variables. Their equality yields \(\mathcal L(S)=\mathcal L(Y)\) through characteristic uniqueness. The next lemma proves the required almost-everywhere correspondence for every actual Verlet iterate by induction from the lifted kick and drift identities. Thus the two convergence arguments refer to precisely the same implemented approximants.



\noindent Lean declarations: \texttt{measure\_map\_lpPhase\_eq\_pointwise\_limit}, \texttt{lp\_verlet\_iterate\_coe}.

\section{Identify the actual coupled representatives almost everywhere}


Marginal equality alone is insufficient when the local numerical error must be coupled with the same initial state and thermostat inputs. This stronger result identifies the actual representatives almost everywhere. Each component of an L\ensuremath{{}^{2}}-convergent phase sequence converges in measure. The known almost-everywhere equality transfers this convergence to the same pointwise Verlet approximants. Their almost-everywhere convergence also implies convergence in measure, because the source is a probability measure. Uniqueness of the limit in measure therefore gives equality almost everywhere of each component, hence of the phase pair.



Apply this argument to the actual L\ensuremath{{}^{2}} Hamiltonian flow and its actual Verlet approximants. The result states that its representative at time t is almost everywhere the finite-dimensional Hamiltonian flow evaluated at the same representative of the initial state. Thus adjoining any initial state or thermostat variable preserves the coupling used in the local comparison; no equality of marginal laws is substituted for a joint-law argument.



\noindent Lean declarations: \texttt{lpPhaseCoe\_eq\_pointwise\_limit}, \texttt{lpHamiltonianFlow\_coe}.

\section{Construct the true L\ensuremath{{}^{2}} Hamiltonian integral flow and identify its law}


The L\ensuremath{{}^{2}} force is the actual Nemytskii lift \texttt{lpDrift}; its global Lipschitz bound permits the existing Picard construction with zero friction and zero noise. Its specification gives a continuous global path, the prescribed initial value, and the actual Hamiltonian integral equation on the finite horizon. Verlet convergence holds in this Banach space. The same vector Verlet maps converge pointwise to the selected finite-dimensional Hamiltonian flow. The proved almost-everywhere pointwise identification and the composition rule for pushforwards give the exact law identity \(\mathcal L(A_t)=(H_t)_\#\mathcal L(A_0)\). This does not infer a pathwise identification from an unproved interchange of representatives and time integrals.



\noindent Lean declarations: \texttt{lpHamiltonianFlow}, \texttt{lpHamiltonianFlow\_spec}, \texttt{lpHamiltonianFlow\_law}.

\section{Derive stationarity and both moment envelopes from the actual target}


If the initial phase law is the canonical independent Gaussian-times-Gibbs target, the exact Hamiltonian invariance theorem and the derived output-law identity prove that the actual L\ensuremath{{}^{2}} flow has that target law at every time. Projecting the actual integral equation with the continuous linear coordinate maps gives the velocity and kinematic component equations; commuting those maps with Bochner integrals is justified by continuity and interval integrability. Projecting the target law gives standard Gaussian velocity and normalized Gibbs position. The exact Gaussian coordinate calculation therefore yields \(\|V_t\|_{L^2}=\sqrt d\), and the already proved Gibbs score estimate gives \(\|\nabla f(X_t)\|_{L^2}\le\sqrt{Md}\). These are conclusions from the actual law, not supplied moment hypotheses.



\noindent Lean declarations: \texttt{lpHamiltonianFlow\_preserves\_target}, \texttt{lpHamiltonianFlow\_component\_equations}, \texttt{lpHamiltonianFlow\_target\_moments}.

\section{Apply the local OBABO comparison to this constructed stationary reference}


Start the genuine L\ensuremath{{}^{2}} Hamiltonian path at the actual first thermostat half-step and assume its phase law is the canonical target, as supplied by thermostat invariance. Its continuity, initial value, velocity equation, kinematic equation, and uniform moments have all been proved above. The local O–Verlet–O versus O–Hamiltonian–O theorem therefore applies with \(C=\sqrt d\) and \(G=\sqrt{Md}\). It yields the actual weighted error bound

\[2h^2\left(\sqrt{Md}+\sqrt M\sqrt d+\frac{h}{4}\sqrt M\sqrt{Md}\right).\]

The two actual thermostat noises are retained, their common final contribution cancels, and the damping needs only \(|\delta|\le1\). This provides the numerical local budget for the global sampling-bias estimate.



\noindent Lean declarations: \texttt{lp\_obabo\_target\_local\_bound}.

\end{document}
